%! Polynomial Rates of Change and Characteristics — Unit 3, Lesson 3.1 — Guided Notes (ANSWER KEY)
% THE in-class document, in the MAIN IDEAS / NOTES shape (LESSON_SHAPE.md §1): the vocabulary
% box, the hook, then ONE two-column guidednotes table. Each row is one idea — a label on the
% left; on the right one or two COMPLETE printed sentences, ONE large pre-drawn display, then a
% two-across grid of numbered problems with real writing room. Problems run 1..18.
% Rates, extrema and concavity were taught on general graphs in 2.5–2.6; here they are RECALLED
% and applied to one polynomial, p(x) = x^3 - 6x^2 + 9x + 2, which is the spine of rows 1–4.
% The NEW instruction is successive differences (row 1) and extrema for polynomials (row 4).
%   I DO   : the teacher reads the definitions, marks up the displays, works problems 1 and 5,
%            then 9 and 11 at the head of their rows.
%   WE DO  : the class works the rest of each grid, every one cold-called.
%   YOU DO : the last row — problems 15–18, alone, for 15 minutes while the teacher circulates.
% TRAP (row 1, problem 4): successive differences on a table whose inputs are NOT equally spaced.
% THE CRUX (row 4, problems 12–13): "a relative maximum is the highest point of the graph" — the
% cubic p has a relative max of 6 but p(5) = 22 (odd degree: no absolute max), and the quartic q
% has two valleys of which only the lower is the absolute minimum. Its You Do transfer is
% problem 18 (a stock whose relative max $25 is not its six-month high, $41 at the endpoint).
% No sketching anywhere: every graph is pre-drawn; students read, label, and compute.
%
% PAGE LOCKSTEP with notes_key/: the vocabulary write lines -> \ansline, the hook's write
% lines -> \ansline, and the answer argument of every \pcell and \labelbox filled. Nothing
% else differs. Inside a cell `\\` ENDS THE ROW — break lines with \par. A row cannot break
% across pages: the figure and EACH grid row are separate sub-rows (`\\` then `& ...`).
\documentclass[10pt]{article}
\usepackage{precalculus-article}
\usepackage{precalculus-key}
\newcommand{\TallMath}[1]{$\displaystyle #1\rule[-1.4em]{0pt}{3.2em}$}
% p(x) = x^3 - 6x^2 + 9x + 2: rel max (1, 6), inflection (2, 4), rel min (3, 2), y-int (0, 2).
\newcommand{\pcurve}{plot[domain=-0.42:4.55, samples=80, smooth]
  (\x, {\x*\x*\x - 6*\x*\x + 9*\x + 2})}
\newcommand{\paxes}{%
  \draw[linegray, very thin, xstep=1, ystep=2] (-1,-4) grid (5,14);
  \draw[->, thick] (-1,0) -- (5.5,0) node[below right] {\small $x$};
  \draw[->, thick] (0,-4) -- (0,14.8) node[above] {\small $p(x)$};
  \foreach \x in {-1,1,2,3,4,5} \node[below, font=\scriptsize] at (\x,0) {\x};
  \foreach \y in {4,8,12} \node[left, font=\scriptsize] at (0,\y) {\y};}

\begin{document}

\pageheader{Unit 3, Lesson 3.1}{Guided Notes --- Answer Key}

\vspace{0.12in}

\boxguard
\begin{vocabbox}
Fill in each term as we define it together.
\par\vspace{2pt}
\noindent\textbf{\textcolor{plum}{Successive differences:}}\\[1pt]
\ansline{Subtract each output from the next in an equal-step table; then repeat on those differences.}\\[3pt]
\noindent\textbf{\textcolor{plum}{Average rate of change on $[a, b]$:}}\\[1pt]
\ansline{$\frac{f(b) - f(a)}{b - a}$: change in output over change in input --- the slope of the secant line.}\\[3pt]
\noindent\textbf{\textcolor{plum}{Difference quotient:}}\\[1pt]
\ansline{$\frac{f(x + h) - f(x)}{h}$, $h \ne 0$: the average rate of change on $[x, x + h]$.}\\[3pt]
\noindent\textbf{\textcolor{plum}{Relative (local) maximum / minimum:}}\\[1pt]
\ansline{A peak / a valley: higher / lower than every nearby output, where the graph turns.}\\[3pt]
\noindent\textbf{\textcolor{plum}{Absolute (global) maximum / minimum:}}\\[1pt]
\ansline{The greatest / least output on the whole domain --- a graph may not have one.}\\[3pt]
\noindent\textbf{\textcolor{plum}{Point of inflection:}}\\[1pt]
\ansline{Where the concavity changes --- the rate of change switches between rising and falling.}
\end{vocabbox}

\vspace{0.1in}

\boxguard
\begin{hookbox}
A roller coaster's first hill is 40 meters tall. Jay says, ``So 40 meters is the highest the
coaster ever goes.''

Is Jay right? What would you need to see to be sure?
\ansline{Not necessarily: the first hill is only the highest point \emph{nearby}. A later hill could be taller --- you need the heights over the whole ride.}
\end{hookbox}

\small
\begin{guidednotes}
% ── ROW 1: successive differences — the degree from an equal-step table (NEW) ─────────────
\mainidea[What degree is it?]{Successive Differences} &
In a table whose inputs go up by \textbf{equal steps}, subtract each output from the next to
get the \textbf{first differences}; subtract those to get the second differences, and so on.
If the $n$th differences are the first ones that are all the same, the data fit a polynomial
of \textbf{degree $n$}: constant first differences mean linear, second mean quadratic, third
mean cubic. \textbf{If the input steps are not equal, the differences say nothing.}
\\
& \begin{center}
$p(x) = x^3 - 6x^2 + 9x + 2$\par\vspace{4pt}
\renewcommand{\arraystretch}{1.25}
\setlength{\tabcolsep}{4pt}
\begin{tabular}{r|*{11}{c}}
$x$            & 0 &   & 1 &    & 2 &    & 3 &   & 4 &    & 5  \\
$p(x)$         & 2 &   & 6 &    & 4 &    & 2 &   & 6 &    & 22 \\ \hline
1st differences &   & 4 &   & $-2$ &   & $-2$ &   & 4 &   & 16 &    \\
2nd differences &   &   & $-6$ &  & 0 &    & 6 &   & 12 &   &    \\
3rd differences &   &   &   & 6  &   & 6  &   & 6 &   &    &    \\
\end{tabular}
\end{center}
\vspace{2pt}
The \labelbox{1.4cm}{third} differences are constant, so $p$ has degree \labelbox{1.0cm}{3}
\\
& \notesprompt{Check the input steps first. Then find the differences and the degree.}
\begin{probgrid}{|Y|Y|}
\pcell{1}{\renewcommand{\arraystretch}{1.1}\begin{tabular}{|c|c|c|c|c|c|}\hline $x$ & 0 & 1 & 2 & 3 & 4 \\ \hline $f(x)$ & 1 & 3 & 9 & 19 & 33 \\ \hline\end{tabular}\par What degree polynomial fits $f$?}{2.2cm}{First differences 2, 6, 10, 14; second 4, 4, 4. Constant second differences: degree 2.} &
\pcell{2}{\renewcommand{\arraystretch}{1.1}\begin{tabular}{|c|c|c|c|c|c|}\hline $x$ & 0 & 1 & 2 & 3 & 4 \\ \hline $g(x)$ & 7 & 4 & 1 & $-2$ & $-5$ \\ \hline\end{tabular}\par What degree polynomial fits $g$?}{2.2cm}{First differences $-3, -3, -3, -3$ are already constant: degree 1 (linear).} \\ \hline
\end{probgrid}
\\
& \begin{probgrid}*{|Y|Y|}
\pcell{3}{\renewcommand{\arraystretch}{1.1}\begin{tabular}{|c|c|c|c|c|c|}\hline $x$ & 0 & 1 & 2 & 3 & 4 \\ \hline $k(x)$ & 1 & 2 & 9 & 28 & 65 \\ \hline\end{tabular}\par What degree polynomial fits $k$?}{2.2cm}{First 1, 7, 19, 37; second 6, 12, 18; third 6, 6. Constant third differences: degree 3.} &
\pcell{4}{\renewcommand{\arraystretch}{1.1}\begin{tabular}{|c|c|c|c|c|c|}\hline $x$ & 0 & 1 & 3 & 4 & 6 \\ \hline $m(x)$ & 1 & 3 & 7 & 9 & 13 \\ \hline\end{tabular}\par Kim says, ``The first differences are 2, 4, 2, 4 --- not constant, so $m$ is not linear.'' What did Kim miss?}{2.2cm}{The $x$-steps are 1, 2, 1, 2 --- not equal. As rates: $\frac{2}{1}, \frac{4}{2}, \frac{2}{1}, \frac{4}{2}$ all equal 2, so $m$ is linear.} \\ \hline
\end{probgrid}
\\ \hline
% ── ROW 2: average rate of change, the secant line, the difference quotient (recall 2.6) ──
\mainidea[How fast, on average]{Rates of Change} &
The \textbf{average rate of change} of $p$ on $[a, b]$ is $\dfrac{p(b) - p(a)}{b - a}$, the
slope of the \textbf{secant line} through $(a, p(a))$ and $(b, p(b))$. Each first difference
in row 1 is an average rate of change over an interval 1 unit wide. The \textbf{difference
quotient} $\dfrac{f(x + h) - f(x)}{h}$ is the average rate of change on $[x, x + h]$: build
$f(x + h)$ by replacing \emph{every} $x$ with $(x + h)$.
\\
& \begin{minipage}[c]{0.62\linewidth}\centering
\begin{tikzpicture}[x=0.85cm, y=0.17cm, >=stealth]
  \paxes
  \draw[very thick, plum, <->] \pcurve;
  % secant through (0, 2) and (1, 6), extended a little
  \draw[very thick, goldacc, dashed] (-0.45,0.2) -- (1.6,8.4);
  \draw[thick, redacc] (0,2) -- (1,2) -- (1,6);
  \node[below, font=\scriptsize\bfseries\color{redacc}] at (0.5,2) {run};
  \node[right, font=\scriptsize\bfseries\color{redacc}] at (1,3.6) {rise};
  \foreach \p in {(0,2),(1,6),(2,4),(3,2),(4,6)} \fill[plum] \p circle (2.3pt);
  \node[font=\scriptsize\itshape\color{goldacc}, left] at (1.5,9.6) {secant line};
\end{tikzpicture}
\end{minipage}\hfill
\begin{minipage}[c]{0.36\linewidth}
Run: \labelbox{1.0cm}{1}\par\vspace{4pt}
Rise: \labelbox{1.0cm}{4}\par\vspace{4pt}
Slope of the secant line:\par \labelbox{1.0cm}{4}
\end{minipage}
\\
& \notesprompt{Use $p(x) = x^3 - 6x^2 + 9x + 2$ and its table in row 1.}
\begin{probgrid}{|Y|Y|}
\pcell{5}{The average rate of change of $p$ on $[0, 2]$.}{2.0cm}{$\frac{p(2) - p(0)}{2 - 0} = \frac{4 - 2}{2} = 1$.} &
\pcell{6}{The average rate of change of $p$ on $[1, 3]$. What does its sign say about $p$ on
that interval?}{2.0cm}{$\frac{p(3) - p(1)}{3 - 1} = \frac{2 - 6}{2} = -2$. Negative: on average $p$ decreases on $[1, 3]$.} \\ \hline
\end{probgrid}
\\
& \begin{probgrid}*{|Y|Y|}
\pcell{7}{$f(x) = x^2 + 4x$. Find $f(x + h)$, then simplify the difference quotient.}{2.6cm}{$f(x + h) = x^2 + 2xh + h^2 + 4x + 4h$. \par Subtract $f(x)$: $2xh + h^2 + 4h$. \par Divide by $h$: $2x + h + 4$.} &
\pcell{8}{Evaluate your answer to problem 7 at $x = 1$, $h = 2$. Check it against the average
rate of change of $f$ on $[1, 3]$.}{2.6cm}{$2(1) + 2 + 4 = 8$. \par Check: $f(1) = 5$, $f(3) = 21$, so $\frac{21 - 5}{3 - 1} = 8$. They match.} \\ \hline
\end{probgrid}
\\ \hline
% ── ROW 3: reading the graph — y-intercept, direction, concavity, inflection (recall 2.5) ─
\mainidea[Reading the graph]{Shape of $p$} &
The \textbf{$y$-intercept} of a polynomial is $(0, p(0))$, and $p(0)$ is its constant term.
Intervals of increasing and decreasing are open intervals of \textbf{$x$-values}. Where the
rate of change is increasing the graph is \textbf{concave up} ($\cup$); where it is decreasing,
\textbf{concave down} ($\cap$). A \textbf{point of inflection} is where the rate of change
switches between increasing and decreasing --- where the concavity changes.
\\
& \begin{minipage}[c]{0.62\linewidth}\centering
\begin{tikzpicture}[x=0.85cm, y=0.17cm, >=stealth]
  \paxes
  \draw[very thick, plum, <->] \pcurve;
  \foreach \p in {(1,6),(3,2),(4,6)} \fill[plum] \p circle (2.3pt);
  \draw[redacc, very thick] (0,2) circle (3.5pt);
  \draw[goldacc, very thick] (2,4) circle (5pt);
  \node[font=\Large\bfseries\color{goldacc}] at (1,10.2) {$\cap$};
  \node[font=\Large\bfseries\color{goldacc}] at (3.75,-2.4) {$\cup$};
\end{tikzpicture}
\end{minipage}\hfill
\begin{minipage}[c]{0.36\linewidth}
The red point is the $y$-intercept: \labelbox{1.4cm}{$(0, 2)$}\par\vspace{6pt}
The gold point $(2, 4)$ is a\par \labelbox{3.2cm}{point of inflection}
\end{minipage}
\\
& \notesprompt{Intervals in $x$-values, with parentheses. The ends of the graph keep going.}
\begin{probgrid}{|Y|Y|}
\pcell{9}{Find $p(0)$ from the formula. Then give the intervals where $p$ is increasing and
where it is decreasing.}{2.2cm}{$p(0) = 2$, the constant term: $(0, 2)$. Increasing on $(-\infty, 1)$ and $(3, \infty)$; decreasing on $(1, 3)$.} &
\pcell{10}{Where is $p$ concave down? Concave up? Then check with the first differences
4, $-2$, $-2$, 4, 16: where do the rates fall, and where do they rise?}{2.2cm}{Concave down on $(-\infty, 2)$; concave up on $(2, \infty)$. The rates fall from 4 to $-2$, then rise to 4 and 16: the switch is at $x = 2$.} \\ \hline
\end{probgrid}
\\ \hline
% ── ROW 4: THE CRUX — relative is not absolute ───────────────────────────────────────────
\mainidea[Highest nearby, or highest of all?]{Relative vs.\ Absolute} &
A \textbf{relative maximum} is higher than the points \emph{near} it --- not necessarily the
highest point of the graph. An \textbf{odd-degree} polynomial has one end rising and one
falling forever, so it has no absolute maximum or minimum. An \textbf{even-degree} polynomial
has both ends going the same way, so it has an absolute minimum (the lowest valley) or an
absolute maximum (the highest peak). On a restricted domain, check the \textbf{endpoints} too.
\\
& \begin{minipage}[c]{0.49\linewidth}\centering
\begin{tikzpicture}[x=0.62cm, y=0.14cm, >=stealth]
  \paxes
  \draw[very thick, plum, <->] \pcurve;
  \foreach \p in {(1,6),(3,2)} \fill[plum] \p circle (2.2pt);
\end{tikzpicture}

{\footnotesize $p(x) = x^3 - 6x^2 + 9x + 2$, degree 3}
\end{minipage}\hfill
\begin{minipage}[c]{0.49\linewidth}\centering
\begin{tikzpicture}[x=0.52cm, y=0.2cm, >=stealth]
  \draw[linegray, very thin, step=1] (-4,-6) grid (5,7);
  \draw[->, thick] (-4,0) -- (5.5,0) node[below right] {\small $x$};
  \draw[->, thick] (0,-6) -- (0,7.6) node[above] {\small $q(x)$};
  \foreach \x in {-3,-2,-1,1,2,3,4} \node[below, font=\scriptsize] at (\x,0) {\x};
  \foreach \y in {-4,-2,2,4,6} \node[left, font=\scriptsize] at (0,\y) {\y};
  \draw[very thick, plum, <->] (-3.05,6.8) .. controls (-2.55,1.2) and (-2.35,-2) .. (-2,-2)
     .. controls (-1.2,-2) and (-0.8,3) .. (0,3)
     .. controls (1.2,3) and (1.9,-5) .. (3,-5)
     .. controls (3.7,-5) and (4.05,1.5) .. (4.3,6.8);
  \foreach \p in {(-2,-2),(0,3),(3,-5)} \fill[plum] \p circle (2.2pt);
\end{tikzpicture}

{\footnotesize $q$, a polynomial of degree 4}
\end{minipage}
\\
& \notesprompt{Give each extremum as a value ($y$) and where it occurs ($x$).}
\begin{probgrid}{|Y|Y|}
\pcell{11}{The relative maximum and the relative minimum of $p$.}{2.0cm}{Relative max: 6, at $x = 1$. \par Relative min: 2, at $x = 3$.} &
\pcell{12}{Kai says, ``The relative maximum of $p$ is 6, so 6 is the highest $p$ ever goes.''
Find $p(5)$. Is Kai right? Does $p$ have an absolute maximum?}{2.2cm}{$p(5) = 125 - 150 + 45 + 2 = 22$. Kai is wrong: $22 > 6$. The right end rises forever, so $p$ has no absolute maximum.} \\ \hline
\end{probgrid}
\\
& \begin{probgrid}*{|Y|Y|}
\pcell{13}{Give both relative minima of $q$. Which one is the absolute minimum? Does $q$ have
an absolute maximum? Why or why not?}{2.4cm}{Relative minima: $-2$ at $x = -2$, and $-5$ at $x = 3$. Absolute min: $-5$, at $x = 3$. No absolute max --- both ends rise forever.} &
\pcell{14}{Suppose $p$ is used only for $0 \le x \le 5$. Now what is the absolute maximum of
$p$, and where does it occur? Is it a turning point?}{2.4cm}{22, at $x = 5$ --- an endpoint, not a turning point. The relative max, 6, is not the absolute max here.} \\ \hline
\end{probgrid}
\\ \hline
% ── LAST ROW: YOU DO — alone, while the teacher circulates (15 min) ─────────────────────
\mainidea[You do, alone]{A New Stock} &
The price of one share of a new company's stock is $S(t)$ dollars, $t$ months after it went on
sale, for $0 \le t \le 6$. $S$ is a polynomial.
\\
& \begin{minipage}[c]{0.4\linewidth}\centering
\renewcommand{\arraystretch}{1.1}
\begin{tabular}{|c|c|}
\hline
$t$ (months) & $S(t)$ (\$) \\ \hline
0 & 5  \\
1 & 21 \\
2 & 25 \\
3 & 23 \\
4 & 21 \\
5 & 25 \\
6 & 41 \\ \hline
\end{tabular}
\end{minipage}\hfill
\begin{minipage}[c]{0.58\linewidth}\centering
\begin{tikzpicture}[x=0.85cm, y=0.065cm, >=stealth]
  \draw[linegray, very thin, xstep=1, ystep=5] (0,0) grid (6,45);
  \draw[->, thick] (0,0) -- (6.5,0) node[below right] {\small $t$};
  \draw[->, thick] (0,0) -- (0,48) node[above] {\small $S(t)$};
  \foreach \x in {1,...,6} \node[below, font=\scriptsize] at (\x,0) {\x};
  \foreach \y in {10,20,30,40} \node[left, font=\scriptsize] at (0,\y) {\y};
  \draw[very thick, plum] plot[domain=0:6, samples=80, smooth]
       (\x, {\x*\x*\x - 9*\x*\x + 24*\x + 5});
  \fill[plum] (0,5) circle (2.4pt);
  \fill[plum] (6,41) circle (2.4pt);
\end{tikzpicture}
\end{minipage}
\\
& \notesprompt{On your own. Show your work. Give intervals in months, prices in dollars.}
\begin{probgrid}{|Y|Y|}
\pcell{15}{Use the outputs for $t = 0$ to $4$ to find successive differences. What degree is
$S$? How do you know?}{2.4cm}{First 16, 4, $-2$, $-2$; second $-12$, $-6$, 0; third 6, 6. Constant third differences: degree 3.} &
\pcell{16}{The average rate of change of $S$ on $[0, 2]$ and on $[2, 4]$, with units. Say what
each one means for the stock.}{2.4cm}{$[0, 2]$: $\frac{25 - 5}{2} = 10$ dollars per month --- up \$10 a month on average. $[2, 4]$: $\frac{21 - 25}{2} = -2$ dollars per month --- down \$2 a month.} \\ \hline
\end{probgrid}
\\
& \begin{probgrid}*{|Y|Y|}
\pcell{17}{Where is $S$ increasing? Decreasing? Where is it concave down, concave up, and what
is the point of inflection?}{2.4cm}{Increasing on $(0, 2)$ and $(4, 6)$; decreasing on $(2, 4)$. Concave down on $(0, 3)$, up on $(3, 6)$. Inflection point $(3, 23)$.} &
\pcell{18}{An investor says, ``The stock peaked at \$25 in month 2, so \$25 was the most it was
worth in these six months.'' Is she right? Give the absolute maximum and minimum of $S$ and
when each occurs.}{2.4cm}{No --- \$25 is only a relative max. Absolute max \$41, at $t = 6$ (the end); absolute min \$5, at $t = 0$ (the start).} \\ \hline
\end{probgrid}
\\ \hline
\end{guidednotes}

% ── THE NOTES END AT THE YOU DO ROW; AP Practice follows; homework is DeltaMath ──────

\end{document}
